%---------------------------------------------------------------------
% Chapter 4: Agent Design and Architecture (template sections; see GUIDANCE.md)
% Sources: FACTS.md sections 3-9 and 24; core/investigator.py, core/verifier.py, core/llm.py,
% core/config.py, core/supervisor.py, core/ingest.py, core/db.py, server/app.py (control logic
% identical at 25fb3f8 and HEAD); docs/trace/run1_trace.jsonl (walkthrough, invented record).
%---------------------------------------------------------------------
\chapter{Agent Design and Architecture}
\label{ch:design}

This chapter presents the design that meets the requirements of Chapter~\ref{ch:req}. It shows
the architecture, describes each component and its failure behaviour, names the agentic patterns
with the alternatives that were rejected, and then follows the control flow, the agent state and
one input through the system. Code and prompts are left to Chapter~\ref{ch:impl}.

\section{System Architecture}
\label{sec:design-system-architecture}

Figure~\ref{fig:arch} shows the architecture as built. It has two lanes that run on different
clocks. In the per-pack lane a person waits for an answer about one pack: a client sends photos
or typed details to a local API, perception tools read the pack, the investigator plans and runs
tool calls over the alert index, and a verifier checks the drafted verdict before it is shown. The
background lane runs after each data sync and never involves the person: it fetches the CDSCO lists,
ingests them into the alert index and re-checks every saved pack. Both lanes run on the user's
machine. The only component that may leave the machine is the portal sync, and the only optional
component is the local language model (dashed in the figure), so a check completes with the model,
the tracing server and the network all switched off (NFR1).

\begin{figure}[!htbp]
\centering\setstretch{1}
\begin{tikzpicture}[
  font=\footnotesize,
  box/.style={draw=bmsblue,fill=guidebg,rounded corners=2pt,minimum width=2.6cm,text width=2.4cm,minimum height=0.95cm,align=center},
  agent/.style={box,line width=0.9pt},
  tool/.style={draw=ForestGreen,fill=dobg,minimum width=2.6cm,text width=2.4cm,minimum height=0.95cm,align=center},
  opt/.style={tool,dashed},
  store/.style={draw=gray,fill=gray!10,cylinder,shape border rotate=90,aspect=0.18,minimum width=2.6cm,text width=2.3cm,minimum height=1.15cm,align=center},
  human/.style={draw=BrickRed,fill=dontbg,rounded corners=6pt,minimum width=2.6cm,text width=2.4cm,minimum height=0.95cm,align=center},
  ext/.style={draw=black!60,fill=white,minimum width=2.6cm,text width=2.4cm,minimum height=0.95cm,align=center,double},
  arr/.style={-{Stealth[length=2mm]},thick},
  lab/.style={font=\footnotesize,align=center,inner sep=3pt}]
% columns 1.3 / 7.45 / 13.6 ; rows below
\node[human] (person) at (7.45,0) {Person holding the pack};
\node[box]   (clients) at (7.45,-1.85) {Clients: web, iOS, command line};
\node[box]   (api) at (7.45,-3.9) {Local API (FastAPI, SSE)};
\node[tool]  (perc) at (1.3,-3.9) {Perception tools: QR, OCR, re-read};
\node[store] (state) at (13.6,-3.9) {Check state, traces, memory};
\node[agent] (inv) at (7.45,-6.8) {\textbf{Investigator} planner and tool loop};
\node[opt]   (llm) at (1.3,-6.8) {Local LLM (optional) qwen2.5:3b};
\node[store] (index) at (13.6,-6.8) {Alert index (SQLite)};
\node[agent] (ver) at (1.3,-8.9) {\textbf{Verifier} rules R1--R10};
\node[store] (raw) at (1.3,-11.45) {Raw snapshots and metadata};
\node[box]   (recheck) at (7.45,-11.45) {Cabinet\\re-check\\(no questions)};
\node[box]   (ingest) at (13.6,-11.45) {Ingest: clean, validate, quarantine};
\node[ext]   (portal) at (1.3,-13.45) {CDSCO SUGAM portal};
\node[box]   (sync) at (7.45,-13.45) {Portal sync\\(1.5\,s gap, TLS)};
% lanes
\begin{scope}[on background layer]
\draw[dashed,draw=bmsblue!60,fill=bmsblue!3,rounded corners=4pt] (-0.15,0.75) rectangle (15.05,-9.95);
\draw[dashed,draw=gray!70,fill=gray!4,rounded corners=4pt] (-0.15,-10.25) rectangle (15.05,-14.45);
\end{scope}
\node[anchor=north west,lab,align=left,text=bmsblue] at (-0.05,0.65) {\textbf{Per-pack lane}\\the person waits};
\node[anchor=south east,lab,align=right,text=black!70] at (14.95,-14.35) {\textbf{Background lane}\\(supervisor)};
% per-pack arrows
\draw[arr] ([xshift=-5mm]person.south) -- node[lab,left,xshift=-1.5mm]{photos, typed\\details, answers} ([xshift=-5mm]clients.north);
\draw[arr] ([xshift=5mm]clients.north) -- node[lab,right,xshift=1.5mm]{questions,\\verdict card} ([xshift=5mm]person.south);
\draw[arr] ([xshift=-5mm]clients.south) -- node[lab,left,xshift=-1.5mm]{HTTP request:\\images, answer} ([xshift=-5mm]api.north);
\draw[arr] ([xshift=5mm]api.north) -- node[lab,right,xshift=1.5mm]{SSE: steps,\\question, verdict} ([xshift=5mm]clients.south);
\draw[arr] ([yshift=2mm]api.west) -- node[lab,above,pos=0.72]{image} ([yshift=2mm]perc.east);
\draw[arr] ([yshift=-2mm]perc.east) -- node[lab,below]{fields, provenance} ([yshift=-2mm]api.west);
\draw[arr,{Stealth[length=2mm]}-{Stealth[length=2mm]}] ([yshift=-2mm]api.east) -- node[lab,above]{check state, trace} node[lab,below]{prior, corrections} ([yshift=-2mm]state.west);
\draw[arr] ([xshift=-5mm]api.south) -- node[lab,left]{fields, answer,\\saved state} ([xshift=-5mm]inv.north);
\draw[arr] ([xshift=5mm]inv.north) -- node[lab,right]{snapshot: question\\or result} ([xshift=5mm]api.south);
\draw[arr,dashed] ([yshift=2mm]inv.west) -- node[lab,above]{evidence,\\allowed actions} ([yshift=2mm]llm.east);
\draw[arr,dashed] ([yshift=-2mm]llm.east) -- node[lab,below]{action, thought} ([yshift=-2mm]inv.west);
\draw[arr] (index.west) -- node[lab,below]{alert rows for\\the batch keys} (inv.east);
\draw[arr] ([xshift=-9mm]inv.south) |- node[lab,above,pos=0.75]{draft verdict} ([yshift=2mm]ver.east);
\draw[arr] ([yshift=-2mm]ver.east) -| node[lab,below,pos=0.25]{approved, or \CANTV} ([xshift=-5mm]inv.south);
% background arrows
\draw[arr] (portal.east) -- node[lab,below]{monthly JSON\\(latin-1)} (sync.west);
\draw[arr] (sync.north west) -- node[lab,below left=-1pt,pos=0.55]{raw files, SHA-256} (raw.south east);
\draw[arr] (sync.north east) -- node[lab,below right=-1pt,pos=0.55]{new snapshots} (ingest.south west);
\draw[arr] (ingest.north) -- node[lab,left,pos=0.4]{validated rows,\\atomic commit} (index.south);
\draw[arr] (ingest.west) -- node[lab,above]{after ingest} (recheck.east);
\draw[arr] ([xshift=9mm]recheck.north) -- node[lab,right,pos=0.75]{saved packs,\\same investigator} ([xshift=9mm]inv.south);
\end{tikzpicture}
\caption{Architecture of StripCheck as built. Blue boxes are control components (bold: the agent and
its verifier), green boxes are tools, the dashed box is optional, cylinders are local stores and the
double box is the external data source.}
\label{fig:arch}
\end{figure}

Three decisions shape the figure. The agent is one bounded loop with an independent verifier,
because the question it answers (is this exact batch on a list?) has a fixed set of evidence steps
that a fixed loop can bound, test and trace. The language model sits beside the loop and never on
the path to the person: it chooses between two allowed actions at a branch point, and the rules take
over when it is slow, absent or invalid. The LangGraph graphs of the specification were replaced by
a custom loop whose whole state is one JSON document saved after every run, which keeps the pause
and resume behaviour with less code (decision D18).

\section{Component Description}
\label{sec:design-component-description}

This section describes every box of Figure~\ref{fig:arch}: its responsibility, inputs and outputs,
method and behaviour on failure.

\subsection{Clients and local API}

The web, iOS and command-line clients talk to one local FastAPI server. \texttt{POST /api/checks}
starts a check with one or two images or typed fields, and \texttt{POST /api/\allowbreak checks/<id>/\allowbreak resume}
continues it with an answer. Progress is streamed as server-sent events (\texttt{step},
\texttt{interrupt} for a question, \texttt{final} for the verdict and others) through a queue of 256
messages, and the work continues if the browser disconnects. The API passes the investigator only
the alert rows whose batch key matches the read or typed batch, together with the size and month
coverage of the whole index. A second answer to a running or changed check is refused with HTTP
409, and an unexpected error marks the check as failed. The server binds to \texttt{127.0.0.1};
phone access needs a pairing token.

\subsection{Perception tools}

Perception turns an image into pack fields, each with its source. A barcode reader (zxing-cpp,
then OpenCV) parses a QR or GS1 code; RapidOCR with PP-OCRv6-small models reads the printed lines on
the CPU~\cite{rapidai2021rapidocr,du2020ppocr}; one enhanced re-read is tried when no batch is found;
label-anchored patterns extract batch, product, brand, dates and maker. Each reader runs in a
separate process within a 60\,s deadline per photo. A failed read leaves the field empty, and the
investigator then asks the person to type the batch. If the QR and printed batches disagree, the
printed one is used without confirmed status; if two photos disagree, the person is asked.

\subsection{Planner}

The default planner is a fixed rule order (Section~\ref{sec:design-control-flow-and-agent-state}).
In hybrid mode the local model \texttt{qwen2.5:3b} under Ollama~\cite{ollama2026structured} is
consulted at two branch points: (b) after an empty exact search, choosing \texttt{search\_variants}
or \texttt{draft\_verdict}, and (c) when grading leaves a candidate that needs a person, choosing
\texttt{ask\_user} or \texttt{draft\_verdict}. It receives a JSON evidence object, the allowed
actions and the rule suggestion, and must return an allowed action and a one-sentence thought. On a
timeout (15\,s by default), a connection error, invalid JSON or an action outside the list, the rule
suggestion is used with the source \texttt{rules(fallback)}, and after a connection error the model
is skipped for 30\,s. A thought with a banned phrase is withheld. The model writes no text that the
person reads.

\subsection{Executor and tools}

The investigator runs each chosen action as a tool call and appends the observation to the history
of the check: \texttt{extract\_fields}, \texttt{search\_exact}, \texttt{search\_variants},
\texttt{grade}, \texttt{ask\_user} and \texttt{draft\_verdict} (signatures in
Chapter~\ref{ch:impl}). The tools are deterministic functions over local data. Inputs a tool cannot
use (no data loaded, a batch shorter than 2 or longer than 40 characters, an unreadable date) end the
run with a \CANTV{} draft and a stated reason instead of an exception.

\subsection{Retriever and relevance grader}

Retrieval is keyed on the normalised batch string, since batch codes carry no meaning that a
semantic search could use. \texttt{search\_exact} looks the key up; \texttt{search\_variants}
generates look-alike keys from a table of confusable characters (8 and B, 0 and O, 1 and I and
others) with at most two substitutions and 2{,}000 lookups, ordered by a noisy-channel score from
earlier confirmations~\cite{kernighan1990spelling,brill2000improved}. The variant search also runs
after an exact hit, so that a listed sibling batch is not missed. The grader treats each candidate as
a record-linkage decision~\cite{fellegi1969theory,christen2012survey}: a disagreeing product, brand
or maker rejects it (a batch collision); a substitution needs a character confirmation; a date one
digit off needs a date confirmation; missing product or maker evidence needs more printed details;
agreeing evidence accepts it. At most 20 candidates are graded.

\subsection{Reflector and human approval gate}

The verifier re-checks every draft against the stored state with ten hard rules (R1 to R10,
Table~\ref{tab:verifier}), independently of the drafting code. An alarm needs a corroborated
candidate, a confirmed batch and a citation; \NLISTED{} needs both searches, no pending candidate and
a statement of the lists loaded; every used field needs a source; no text may contain a banned
phrase. A failing draft becomes the \CANTV{} template with the rule ids. Only one draft is sent back
for revision: a model draft made before the variant search, which rule R4 returns once.

The human gate is \texttt{ask\_user}, with five question kinds: type the batch, confirm the batch,
confirm a character, confirm a printed date, or add printed details or a carton photo. An alarm
requires the batch to come from a QR code, from the person or from two agreeing photos, never from
OCR alone. The person types the characters at the look-alike positions instead of pressing a yes
button, because people tend to accept confirmations without
checking~\cite{parasuraman2010complacency,bucinca2021trust}. Every question offers ``Can't tell'',
and a declined question leads to \CANTV.

\subsection{Stores and memory}

The \emph{alert index} is a SQLite table of cleaned rows, each with its source URL, the SHA-256
digest of its snapshot and its row number. The \emph{check state} store keeps one JSON document per
check, written after every run so that a pending question survives a restart, and the same database
holds the cabinet of saved packs, notifications and confirmed character corrections; each check also
writes a JSON Lines trace. The \emph{raw snapshots} are the portal responses with their metadata.
Memory has two parts: the cabinet of saved packs, and a confusion prior of character pairs that
people confirmed on photographed packs. The prior reorders variants and chooses which characters to
ask about; it never removes a question.

\subsection{Background supervisor}

The portal sync fetches the newest three months by default (the history from January 2019 on
request), at least 1.5\,s apart, with TLS verification, three attempts and a 20\,s timeout, and
stores each raw response. Ingest reads the latin-1 JSON, validates the rows, quarantines rows with a
missing product, a placeholder batch or a date in the batch cell, and commits each snapshot in one
transaction; a snapshot whose digest is already stored is skipped, so imports are idempotent. The
cabinet re-check then runs the same investigator on every saved pack with questions disabled and
notifies the person of a new listing or a printed expiry within 30 days. When the portal is
unreachable the existing index stays in use.


\section{Selected Agentic Patterns and Justification}
\label{sec:design-selected-agentic-patterns-and-justification}

Table~\ref{tab:design-patterns} lists the patterns used, where each sits, why it was chosen and
which alternative was rejected. The common reason is the cost of an error: a false alarm on a
sibling batch harms a product that is not listed, and a confident \NLISTED{} can be read as an
assurance about the medicine. Every decision that reaches the person is therefore kept in testable
rules, and the model is used only where a later rule catches a wrong choice, in line with the advice
to prefer simple workflows when the steps are known~\cite{schluntz2024building}.

\begin{table}[!htbp]
\centering
\caption{Agentic patterns used in StripCheck and the alternatives rejected}
\label{tab:design-patterns}
\small\renewcommand{\_}{\textunderscore\allowbreak}
\begin{tabularx}{\linewidth}{@{}>{\raggedright\arraybackslash}p{2.6cm} >{\raggedright\arraybackslash}p{2.2cm} >{\raggedright\arraybackslash}X >{\raggedright\arraybackslash}p{2.9cm}@{}}
\toprule
\textbf{Pattern} & \textbf{Where} & \textbf{Reason} & \textbf{Alternative rejected} \\ \midrule
Bounded tool loop (plan, act, observe)~\cite{yao2023react} & Investigator & Each step is a logged tool call; the loop has a step cap and runs without a model & Free-form ReAct loop driven by the model \\
Model as advisor at branch points~\cite{kambhampati2024llmmodulo} & Planner (b), (c) & A 3B local model chooses only among allowed actions; rules check and replace its choice & Model plans every step \\
Corrective retrieval with a grader & Exact and variant search, grade & A misread batch is expected; an empty exact search triggers a bounded look-alike search, and grading removes collisions & Single exact lookup (the baseline of Chapter~\ref{ch:eval}) \\
Record linkage~\cite{fellegi1969theory} & Grader & A shared batch string is not identity; product, dates and maker must agree & Embedding similarity over alert text \\
External verifier (reflection)~\cite{huang2024cannot} & Verifier & Hard rules re-check the draft against stored state & Model self-critique~\cite{madaan2023selfrefine,shinn2023reflexion} \\
Human-in-the-loop interrupt & \texttt{ask\_user} & Only the person holding the pack can read the printed character & Yes/no confirm button; fully automatic alarm \\
Abstention~\cite{chow1970reject,geifman2017selective} & \CANTV{} & Weak evidence ends in an explicit ``more evidence is needed'' & Forced choice between alarm and \NLISTED \\
Supervisor with re-check & Background lane & Sequential sync, ingest, re-check with one investigator & Multi-agent conversation~\cite{wu2023autogen} \\
\bottomrule
\end{tabularx}
\end{table}

A model-driven ReAct loop was rejected because a 3B model on a laptop adds seconds per decision
and its choices cannot be bounded in advance; Chapter~\ref{ch:exp} later found that the hybrid
planner changed no verdict on the \ResN{} evaluation packs while its 95th-percentile latency was
about five times that of the rule planner. Model self-critique was rejected for the reason given
in Section~\ref{sec:lit-related-technical-work}: without external feedback it does not reliably
correct the model's reasoning.

\section{Control Flow and Agent State}
\label{sec:design-control-flow-and-agent-state}

Figure~\ref{fig:design-flow} shows the decision flow of one run of the investigator. A run starts
from the supplied fields, or resumes from the saved state with an answer. Every path ends in
\texttt{draft\_verdict} and the verifier, or in a question that pauses the check; there is no other
exit. The diamonds are the decision points, and the two branch points where the hybrid planner may
consult the model are marked (b) and (c).

\begin{figure}[!htbp]
\centering\setstretch{1}
\begin{tikzpicture}[
  font=\footnotesize,
  stp/.style={draw=bmsblue,fill=guidebg,rounded corners=2pt,minimum width=3.6cm,text width=3.4cm,minimum height=0.75cm,align=center},
  dec/.style={draw=bmsblue,fill=white,diamond,aspect=3.2,inner sep=0.5pt,align=center},
  hum/.style={draw=BrickRed,fill=dontbg,rounded corners=4pt,text width=2.7cm,minimum height=0.9cm,align=center},
  side/.style={draw=ForestGreen,fill=dobg,rounded corners=2pt,text width=3.0cm,minimum height=0.8cm,align=center},
  bad/.style={draw=black!60,fill=gray!12,rounded corners=2pt,text width=2.9cm,minimum height=0.8cm,align=center},
  arr/.style={-{Stealth[length=2mm]},thick},
  lab/.style={font=\footnotesize,align=center,inner sep=1.5pt}]
\node[stp] (start) at (6.4,0) {Start, or resume with an answer};
\node[dec]  (d1) at (6.4,-1.5) {Inputs usable?};
\node[stp] (ex) at (6.4,-3.0) {\texttt{search\_exact}};
\node[dec]  (d2) at (6.4,-4.5) {Exact hit?};
\node[stp] (va) at (6.4,-6.0) {\texttt{search\_variants}};
\node[stp] (gr) at (6.4,-7.5) {\texttt{grade} candidates};
\node[dec]  (d3) at (6.4,-9.0) {Candidate status?};
\node[dec]  (d4) at (6.4,-10.6) {Gate evidence?};
\node[stp] (dr) at (6.4,-12.1) {\texttt{draft\_verdict} (template)};
\node[dec]  (d5) at (6.4,-13.6) {Verifier passes?};
\node[stp,fill=dobg,draw=ForestGreen] (out) at (6.4,-15.1) {Verdict with citation};
\node[hum]  (ask) at (1.5,-6.75) {\texttt{ask\_user}: pause, state saved};
\node[side] (pb) at (11.6,-4.5) {Branch (b): rules or LLM choose; a draft is sent back once (R4)};
\node[bad]  (cv) at (11.6,-13.6) {\CANTV{}\\template, rule ids};
\node[draw=black!50,fill=white,text width=2.9cm,align=left,inner sep=3pt] (bounds) at (1.5,-13.9)
  {\textbf{Bounds}\\16 steps (R6)\\2 questions\\no question after step 12\\2 substitutions\\2{,}000 variant keys\\20 candidates\\1 send-back\\LLM: 2 tries, 15\,s};
% spine
\draw[arr] (start) -- (d1);
\draw[arr] (d1) -- node[lab,right]{yes} (ex);
\draw[arr] (ex) -- (d2);
\draw[arr] (d2) -- node[lab,right]{yes} (va);
\draw[arr] (va) -- (gr);
\draw[arr] (gr) -- (d3);
\draw[arr] (d3) -- node[lab,right]{accepted} (d4);
\draw[arr] (d4) -- node[lab,right]{yes: \FLAGGED{} or \SFAKE} (dr);
\draw[arr] (dr) -- (d5);
\draw[arr] (d5) -- node[lab,right]{pass} (out);
% branch (b) and send-back
\draw[arr] (d2) -- node[lab,above]{no} (pb);
\draw[arr] (pb.south) |- node[lab,above,pos=0.75,yshift=1.5mm]{\texttt{search\_variants}} (va.east);
% questions
\draw[arr] (d1.west) -| node[lab,above,pos=0.3]{no batch, or\\photos disagree} (ask.north);
\draw[arr] (d3.west) -| node[lab,above,pos=0.3,yshift=1.5mm]{pending: (c)} ([xshift=5mm]ask.south);
\draw[arr] (d4.west) -| node[lab,above,pos=0.3,yshift=1.5mm]{no: \texttt{confirm\_batch}} ([xshift=-5mm]ask.south);
\draw[arr,dashed] (ask.west) -- ++(-0.45,0) |- node[lab,above,pos=0.75,yshift=1.5mm]{answer: resume} (start.west);
% abstention rail
\draw[thick] (d1.east) -- node[lab,above]{no data, invalid length\\or date, step limit: \CANTV} (14.3,-1.5);
\draw[arr] (14.3,-1.5) |- (dr.east);
\draw[thick] (d3.east) -- node[lab,above]{none left: \NLISTED;\\conflict or too many: \CANTV} (14.3,-9.0);
\draw[arr] (d5) -- node[lab,above]{fail} (cv);
\draw[arr] (cv.south) |- (out.east);
\end{tikzpicture}
\caption{Decision flow of one investigator run. Diamonds are decisions, (b) and (c) are the branch
points where the hybrid planner may consult the local model, red is the pause for a question, and
grey is the downgrade by the verifier. Every path reaches the verifier or a question.}
\label{fig:design-flow}
\end{figure}

A pending candidate leads to the matching question: \texttt{confirm\_char} for a substituted
character, \texttt{confirm\_field} for a date one digit off and \texttt{box\_photo} for missing
product or maker evidence. When no question may be asked, or the planner at (c) drafts, the reason
is drafted as \CANTV. After an answer the check resumes from the saved state, applies the answer to
the fields and runs the searches and the grader again. The state is one JSON document per check
(Table~\ref{tab:design-state} in Appendix~\ref{app:supp}), stored by the API after every run and loaded by the resume endpoint,
so a check can wait for an answer across a server restart.

Every loop in the design has an upper bound; Table~\ref{tab:design-bounds} in
Appendix~\ref{app:supp} lists the termination condition of each. The values are constants in \texttt{core/config.py} and the reader,
model and streaming modules, and settings read from the environment are clamped; the verifier
enforces the step and question limits again as rule R6, so a defect in the investigator cannot
produce an unbounded check that is still shown as a result.

\section{Data Flow}
\label{sec:design-data-flow}

This section follows one input end to end. The example is the recorded walkthrough of 8 November
2026 (\texttt{docs/trace/run1\_trace.jsonl}), run with the hybrid planner against a throwaway
database that holds one invented alert record (batch \texttt{LM2408B} for ``Lumera Tablets'' by
``Example Laboratories'', listed as NSQ in September 2026), so that no CDSCO row is copied. The pack
carries batch \texttt{LM2408B}, but the reader returned \texttt{LM24088}: the final B was read as 8.

Perception and the API pass the investigator the fields batch \texttt{LM24088}, product ``Lumera
Tablets'', manufacture 2024-08, expiry 2027-07 and the maker, each with provenance \texttt{ocr} and
batch evidence \texttt{ocr}. The exact search (step~2) returns 0 candidates. At branch point (b) the
model receives the read batch, the 0 exact hits and the allowed actions and returns
\texttt{search\_variants}, which finds one candidate, \texttt{LM2408B}, with the substitution 8 to B
at position~7 (step~3). Grading finds product, dates and maker in agreement and sets the status
\texttt{needs\_confirm\_char} (step~4). At branch point (c) the model returns \texttt{ask\_user}, the
check stores its state and pauses with a \texttt{confirm\_char} question (step~5). The person chooses
\texttt{LM2408B}; the batch is replaced, its evidence becomes \texttt{user\_confirmed}, and the pair
(8, B) is stored in the confusion prior (step~6). The searches and the grader run again on the
confirmed batch and accept the candidate (steps~7 to~9), and the \FLAGGED{} template with list,
month and source passes the verifier (step~10). Table~\ref{tab:app-trace1} in
Appendix~\ref{app:traces} lists the recorded steps.

A plain lookup would have answered \NLISTED{} at step 2. The model chose the same actions as the
rules at both branch points, and the check used 10 of its 16 steps and one of its two questions.
